\section{Main Results}
\label{sec:results}

\subsection{Uninsured Time Deposits}
\label{sec:results:uninsured}

Uninsured time deposits fall after the disclosure, in proportion to the
disclosed charge.  Table~\ref{tab:small_bank_did_time_deposits},
columns (1)--(2), reports the baseline difference-in-differences
estimates.  High-CECL banks lost 0.394 percentage points of assets in
uninsured time deposits after adoption ($p<0.05$), roughly 7 percent of
the sample mean; the continuous specification implies a loss of 0.072
percentage points per percentage point of CECL exposure ($p<0.05$).
Time deposits are the category where the sharpest response should
appear: certificates are held for investment rather than transaction
purposes, they reprice only at maturity, and the rollover decision
belongs to the depositor \citep{calomiris1991role, iyer2016tale,
maechler2006dynamic, egan2017deposit}.  The disclosure arrives while
that decision is on the table.

The rollover mechanism implies the outflow should be concentrated
where certificates were actually coming due.  To check this, I split
the sample at the median of a predetermined measure of maturity
exposure: the share of each bank's uninsured time deposits with
remaining maturity under one year as of 2022Q4.  The estimate in the
high-maturity half is $-0.544$ (SE $0.280$), 2.4 times the $-0.229$
(SE $0.236$) in the low-maturity half: the outflow concentrates where
rollover was feasible.  Two caveats keep this corroboration rather
than an independent test.  With only 213 high-CECL banks split across
the two halves, the formal difference is imprecise (interaction
$-0.21$, SE $0.36$), and the 2022Q4 maturity structure is itself
partly bank-chosen.

\subsection{Insured Time Deposits}
\label{sec:results:insured}

The mirror image appears on the insured side.  Columns (3)--(4) of
Table~\ref{tab:small_bank_did_time_deposits} estimate the same
specification for insured time deposits: high-CECL banks \emph{gained}
0.712 percentage points of assets ($p<0.05$), with a matching
continuous estimate of $+0.070$.  Funding did not simply leave
high-CECL banks; it changed insurance status within the same
instrument class.  The symmetry of the two estimates is the first
indication that the response is a reallocation across the insurance
threshold rather than a general funding outflow.

\subsection{Event-Study Evidence}
\label{sec:results:eventstudy}

Figure~\ref{fig:overlay_composition} overlays the event-study paths of
the two series for a ten-quarter window on either side of adoption
($k=-1$ omitted).  Pre-adoption coefficients for uninsured time
deposits are clustered near zero across all ten lead quarters, and the
uninsured series breaks downward within the first post-adoption quarter
and stays down through the end of the sample.  The insured series rises
after adoption as well, but its movement starts slightly earlier, in
the final pre-adoption quarters; that early start is the first trace of
the bank-side anticipation documented at the end of this section.

Figure~\ref{fig:es_long_balanced} extends the window back to 2018Q1.
Because 92 percent of banks adopt in the same quarter, the most distant
event times are reached only by a handful of off-cycle adopters, so the
two outermost coefficients ($k=-20$ and $k=+11$) pool all event times
at or beyond those points rather than estimating each edge quarter from
a few banks.  Over the five-year pre-period, point estimates for
uninsured time deposits sit near zero with no slope toward the
adoption-date divergence; joint tests of the pre-adoption coefficients
do not reject parallel trends (uninsured $p=0.18$, insured $p=0.37$).
The COVID quarters contribute noise but no differential trend.  The one
visible pre-adoption movement is a mild late upward drift in the
\emph{insured} series (Panel~B).  That drift is the balance-sheet
footprint of the banks' own anticipatory pre-funding: brokered
certificates are booked as time deposits, and banks that knew large
charges were coming began raising insured brokered money before
adoption, bank-side behavior I document directly below.  Because the
reference period is $k=-1$, this pre-positioning is absorbed into the
baseline rather than attributed to the post-adoption response.

\subsection{The Composition Shift: Triple Difference}
\label{sec:results:composition}

As a deliberately conservative check on the level results, the triple
difference of equation~\eqref{eq:triple} compares uninsured to insured
time deposits within the same bank and quarter under bank $\times$
quarter fixed effects.
Table~\ref{tab:small_bank_triple_diff_time_deposits} reports the
estimates.  The binary triple interaction is $-1.252$ (SE $0.340$,
$p<0.01$): relative to their own pre-adoption composition and relative
to low-CECL banks, high-CECL banks experienced a 125-basis-point
widening of the insured-to-uninsured gap as a share of assets after
adoption.  The continuous analog ($-0.177$, $p<0.01$) confirms a
smooth dose-response.  Because the fixed effects absorb any level
shift in total time-deposit demand at each bank in each quarter,
including any bank-specific effect of the 2023 stress, whatever
differential stress exposure high-CECL banks carried cannot explain
why uninsured deposits fell \emph{relative to insured deposits within
the same bank} in proportion to the disclosed adjustment.

One caution on interpretation.  Both arms of the differential are
deposit categories, so the comparison is between two potentially
endogenous outcomes rather than between the treated outcome and an
unaffected control liability \citep{covitz2004market}.  The insured
increase in part reflects active bank substitution into insured
brokered and reciprocal deposits, so the triple difference measures the
joint depositor-plus-bank response at the insurance margin, not a
depositor-only quantity.

\subsection{Discipline and the Capital Buffer}
\label{sec:results:capital}

The disclosed charge should matter most where the buffer absorbing it
is thinnest: a given Day-One adjustment moves a thin-capital bank
closer to distress than a well-capitalized one, so a monitoring
depositor should respond more strongly.  Table~\ref{tab:capital_het}
tests this by interacting the DiD terms with each bank's predetermined
2022Q4 equity-to-assets ratio.  Columns (1)--(3) use the level of
uninsured time deposits; columns (4)--(6) use the within-bank gap
between uninsured and insured time deposits, the outcome equivalent to
the bank-by-quarter triple difference.

The treatment effect is largest at low capital and attenuates as the
buffer grows, in every specification.  In the gap design, each
percentage point of CECL exposure moves the composition by $-0.80$
percentage points of assets at a hypothetical zero-equity bank
($p<0.01$), attenuating by $0.067$ per point of capital ($p<0.01$);
the binary analog is $-4.12$ and $+0.31$.  The implied effect reaches
zero at an equity ratio of roughly 12 to 13 percent, just above the
75th percentile of the sample distribution: at the best-capitalized
quartile of banks, the same disclosed charge draws essentially no
response.  The estimates are unchanged when the time-varying equity
control is dropped (columns 2 and 5), so the interaction is not an
artifact of controlling for contemporaneous capital; the one imprecise
cell is the binary treatment in levels (column 3), which the
composition design sharpens.  This gradient is what monitoring theory
predicts: the disclosure is news about distance to default, and
distance to default depends on the buffer the charge consumes.

\subsection{No Repricing: Two Rate Tests}
\label{sec:results:rates}

Did high-CECL banks defend their uninsured funding by paying more for
it?  Two independent rate datasets say no.

Table~\ref{tab:td_rates} reports DiD estimates for implied deposit
rates computed from Call Report interest expense on time deposits above
and below the \$250{,}000 insurance threshold.  The rate high-CECL
banks pay on large (uninsured) time deposits does not rise relative to
peers after adoption, and the within-bank spread between large- and
small-denomination rates, if anything, narrows slightly.  There is no
evidence of a risk premium extracted by the uninsured depositors who
stayed.

Table~\ref{tab:rw_did} takes the same question to posted retail rates:
a bank-week panel of RateWatch offered APYs covering roughly 3{,}300
sample banks, with the post period split at the date the disclosures
became public.  Posted 12-month CD rates at the \$10{,}000 and
\$100{,}000 tiers and posted savings rates at high-CECL banks show no
differential movement after publication (e.g.\ $-0.058$, SE $0.074$,
for the retail CD), and none during the March 10--April 30, 2023
stress window either ($-0.026$, SE $0.048$).  High-CECL banks never
competed on retail price, before or after the news.

\subsection{Repurchasing Insurance: Where the Funding Goes}
\label{sec:results:destination}

If banks did not defend uninsured balances on price, how did they
replace them?  Table~\ref{tab:brokered} examines the wholesale-funding
margin.  After adoption, high-CECL banks increase insured brokered
deposits by 0.52 percentage points of assets ($p<0.05$) and reciprocal
deposits, the network arrangement that purchases full insurance
coverage for large accounts, by 0.089 percentage points per unit of
continuous CECL exposure ($p<0.10$), while \emph{uninsured} brokered
funding does not move.  Because insured brokered and reciprocal
deposits are two routes to the same destination, insured wholesale
funding, Table~\ref{tab:ins_wholesale} combines them into a single
outcome: high-CECL banks increase the combined position by 0.81
percentage points of assets after adoption ($p<0.05$; continuous
$0.175$, $p<0.05$), a sharper estimate than either component alone.
The margin is intensive rather than extensive: the probability of
holding any reciprocal balance does not move.  High-CECL banks did not
lose funding so much as repurchase it in insured form, the
market-based insurance margin of \citet{kim2024insurance}.  Excluding
the quartile of banks with the largest post-adoption growth in insured
brokered deposits shrinks the triple interaction toward zero
(unreported), indicating that the composition shift is concentrated
precisely among the banks that actively replaced uninsured funding
with insured wholesale money.

The event-time paths make the margin visible.  Insured brokered
deposits (Figure~\ref{fig:es_brokered_ins}) are flat through event
quarter $-4$, begin accumulating around three quarters \emph{before}
adoption, and continue rising through the end of the sample; the
pre-adoption phase is the bank-side anticipation I analyze below.
Reciprocal deposits (Figure~\ref{fig:es_reciprocal}) trend upward
across the whole window as reciprocal networks expanded secularly,
then step up further in the later post-adoption quarters; the
confidence bands are wide, and I treat the DiD coefficients in
Tables~\ref{tab:brokered} and~\ref{tab:ins_wholesale} as the estimates
and the figure as context.

\subsection{Profitability}
\label{sec:results:financial}

The substitution is costly.
Table~\ref{tab:small_bank_did_financial_outcomes} reports DiD estimates
for interest expense, ROA, and NIM.  High-CECL banks paid $0.021$
percentage points more per quarter in interest expense per dollar of
assets ($p<0.01$), approximately 8--9 basis points annualized, or 9
percent of the sample quarterly mean.  ROA fell by $0.038$ percentage
points per quarter ($p<0.01$), an annualized 15-basis-point drag, and
NIM contracted by $0.024$ percentage points per quarter ($p<0.05$).
The joint decline in ROA and NIM is consistent with higher liability
costs not being passed through to borrowers, in line with community
banks facing competitive local loan markets \citep{hannan1988bank}.
Continuous estimates scale proportionally.  Combined with the rate
tests above, the source of the higher funding cost is the substitution
itself: maturing uninsured certificates were replaced with insured
brokered and reciprocal money raised at wholesale rates, not repriced
in place.

\subsection{Bank Anticipation and the Timing of the Response}
\label{sec:results:anticipation}

Banks knew their Day-One numbers before depositors could.  Supervisors
and auditors expected community banks to run CECL in parallel with the
incurred-loss model through 2022, so management had a quantified
estimate of the coming capital charge before adoption; the public
learned it only when 2023Q1 Call Reports were posted in late April and
May 2023.  If the funding reallocation documented above reflects
responses to CECL information rather than to a common stress shock,
bank-controlled margins should move on the private information, before
adoption, and depositor-facing margins only on the public release.
Both predictions hold.

Table~\ref{tab:anticipation} tests the bank side.  Columns (1)--(2)
regress each bank's change in insured brokered deposits over event
quarters $-4$ to $-1$, before any public disclosure, on the size of
its still-undisclosed future adjustment.  The pre-adoption ramp scales
with the future charge: 0.12 percentage points of assets per percentage
point of CECL exposure ($p<0.01$), 0.53 percentage points at high-CECL
banks.  Column (3) shows the same private knowledge in a regulatory
choice: the CECL capital transition election, made on the first Call
Report a bank files under CECL and irrevocable thereafter, allows the
bank to phase the Day-One hit into regulatory capital over three
years, adding back 75 percent of the impact in the first year, 50
percent in the second, and 25 percent in the third, with the full
impact reflected from the fourth year on.  The election is
economically sensible only for a bank expecting a material hit, and
the probability of electing rises steeply with the charge: 34 percent
of high-CECL banks elected relief against 6 percent of the rest.
Column (4) closes the loop: the banks that elected relief are the
banks that pre-funded, with a pre-adoption brokered ramp 0.41
percentage points larger ($p<0.05$).

The depositor side shows the opposite timing.  Uninsured time deposits
are flat through every pre-adoption quarter
(Figure~\ref{fig:overlay_composition}; five-year joint test $p=0.18$)
and break only at adoption, when the information became public.
Posted deposit rates sharpen the contrast at weekly frequency: high-CECL
banks' posted CD and savings rates show no differential movement in the
weeks after the March 2023 failures (Table~\ref{tab:rw_did}), which is
difficult to square with the view that these banks were differentially
caught in the 2023 panic, and no movement after publication either,
consistent with the wholesale, not retail, replacement channel.  A
common stress shock would move bank and depositor margins together in
March 2023.  Instead, each margin moves exactly when its
decision-maker's information arrives: banks in 2022 on private
estimates, depositors at disclosure on the public number.
